% GENERATED FILE. Edit canonical lessons, then run npm run generate:books.
% canonical-source-hash: fnv1a64:158f947e29d72054
% canonical-lessons: TA-C303-kuccal, TA-C303-amali, TA-C303-kuttanerical, TA-C303-tancam, TA-C303-virattu

\chapter{A Crowd Out of Control}
\label{ch:ta-a-crowd-out-of-control}
\hlchaptermodality{303}

\begin{chapteropening}
\textbf{By the end of this chapter:} \emph{I can describe a noisy crowd and what it did.}

Complete the last lesson of chapter 303: I can describe a noisy crowd and what it did.
\end{chapteropening}

\section[\texorpdfstring{\ta{கூச்சல்} (kūccal)}{கூச்சல் (kūccal)}]{\ta{கூச்சல்} (kūccal) — an uproar; an outcry}

\label{lesson:TA-C303-kuccal}



\begin{quote}
Before the new one: say the Tamil for to rise again, then the Tamil for to ravage.
\end{quote}

\subsection*{You'll want to know: \ta{கூச்சல்}}
\textbf{\ta{கூச்சல்}} — \emph{kūccal} — \textquotedblleft{}an uproar; an outcry\textquotedblright{}.

\begin{quote}
\ta{வெளியே} \ta{பெரிய} \ta{கூச்சல்} \ta{கேட்டது}.
\end{quote}

\begin{quote}
\emph{veḷiyē periya kūccal kēṭṭatu.}
\end{quote}

\begin{quote}
A loud uproar was heard outside.
\end{quote}

\ta{கூச்சல்} \ta{போடு}, to make a racket.

\begin{cousinweb}
\textbf{\ta{கூச்சல்}} (\emph{kūccal}) is native Dravidian (DEDR 1868), from the cry \ta{கூ}, as in a call or a shout.
\end{cousinweb}

\begin{grammarlens}[title={What you've built}]
One more word for telling what happened.
\end{grammarlens}

\subsection*{Your turn}
\begin{itemize}
\raggedright
  \item \emph{Say it:} \emph{kūccal}
  \item \emph{Say it:} \emph{kūccal}, once more
  \item \emph{Say it:} \emph{veḷiyē periya kūccal kēṭṭatu.} — A loud uproar was heard outside.
  \item \emph{Recall:} say the Tamil for to rise again, then the Tamil for to ravage, then say \emph{kūccal} again
\end{itemize}

\begin{tcolorbox}[breakable,colback=teal!4,colframe=teal!35!black,title={Before you move on}]
What does \ta{கூச்சல்} mean? (\textquotedblleft{}An uproar; an outcry\textquotedblright{}.) Say it once more.
\end{tcolorbox}
\section[\texorpdfstring{\ta{அமளி} (amaḷi)}{அமளி (amaḷi)}]{\ta{அமளி} (amaḷi) — a commotion}

\label{lesson:TA-C303-amali}



\begin{quote}
Before the new one: say the Tamil for to ravage, then the Tamil for an uproar.
\end{quote}

\subsection*{You'll want to know: \ta{அமளி}}
\textbf{\ta{அமளி}} — \emph{amaḷi} — \textquotedblleft{}a commotion\textquotedblright{}.

\begin{quote}
\ta{கூட்டத்தில்} \ta{திடீரென்று} \ta{அமளி} \ta{தொடங்கியது}.
\end{quote}

\begin{quote}
\emph{kūṭṭattil tiṭīreṉṟu amaḷi toṭaṅkiyatu.}
\end{quote}

\begin{quote}
A commotion suddenly broke out at the meeting.
\end{quote}

\ta{அமளி} \ta{செய்}, to make a commotion.

\begin{cousinweb}
\textbf{\ta{அமளி}} (\emph{amaḷi}) is native Dravidian (DEDR 166): tumult, uproar, bustle.
\end{cousinweb}

\begin{grammarlens}[title={What you've built}]
One more word for telling what happened.
\end{grammarlens}

\subsection*{Your turn}
\begin{itemize}
\raggedright
  \item \emph{Say it:} \emph{amaḷi}
  \item \emph{Say it:} \emph{amaḷi}, once more
  \item \emph{Say it:} \emph{kūṭṭattil tiṭīreṉṟu amaḷi toṭaṅkiyatu.} — A commotion suddenly broke out at the meeting.
  \item \emph{Recall:} say the Tamil for to ravage, then the Tamil for an uproar, then say \emph{amaḷi} again
\end{itemize}

\begin{tcolorbox}[breakable,colback=teal!4,colframe=teal!35!black,title={Before you move on}]
What does \ta{அமளி} mean? (\textquotedblleft{}A commotion\textquotedblright{}.) Say it once more.
\end{tcolorbox}
\section[\texorpdfstring{\ta{கூட்டநெரிசல்} (kūṭṭanerical)}{கூட்டநெரிசல் (kūṭṭanerical)}]{\ta{கூட்டநெரிசல்} (kūṭṭanerical) — a stampede; a crush}

\label{lesson:TA-C303-kuttanerical}



\begin{quote}
Before the new one: say the Tamil for an uproar, then the Tamil for a commotion.
\end{quote}

\subsection*{You'll want to know: \ta{கூட்டநெரிசல்}}
\textbf{\ta{கூட்டநெரிசல்}} — \emph{kūṭṭanerical} — \textquotedblleft{}a stampede; a crush\textquotedblright{}.

\begin{quote}
\ta{கூட்டநெரிசலில்} \ta{நிறைய} \ta{பேர்} \ta{கீழே} \ta{விழுந்தார்கள்}.
\end{quote}

\begin{quote}
\emph{kūṭṭanericalil niṟaiya pēr kīḻē viḻuntārkaḷ.}
\end{quote}

\begin{quote}
Many people fell down in the stampede.
\end{quote}

\ta{கூட்டம்}, crowd + \ta{நெரிசல்}, crush.

\begin{cousinweb}
\textbf{\ta{கூட்டநெரிசல்}} (\emph{kūṭṭanerical}) is a Tamil compound built from \ta{கூட்டம்}, \textquotedblleft{}crowd\textquotedblright{} (DEDR 1882), and \ta{நெரிசல்}, \textquotedblleft{}a crush\textquotedblright{}, from \ta{நெரி}, \textquotedblleft{}to be crushed\textquotedblright{} (DEDR 2927).
\end{cousinweb}

\begin{grammarlens}[title={What you've built}]
One more word for telling what happened.
\end{grammarlens}

\subsection*{Your turn}
\begin{itemize}
\raggedright
  \item \emph{Say it:} \emph{kūṭṭanerical}
  \item \emph{Say it:} \emph{kūṭṭanerical}, once more
  \item \emph{Say it:} \emph{kūṭṭanericalil niṟaiya pēr kīḻē viḻuntārkaḷ.} — Many people fell down in the stampede.
  \item \emph{Recall:} say the Tamil for an uproar, then the Tamil for a commotion, then say \emph{kūṭṭanerical} again
\end{itemize}

\begin{tcolorbox}[breakable,colback=teal!4,colframe=teal!35!black,title={Before you move on}]
What does \ta{கூட்டநெரிசல்} mean? (\textquotedblleft{}A stampede; a crush\textquotedblright{}.) Say it once more.
\end{tcolorbox}
\section[\texorpdfstring{\ta{தஞ்சம்} (tañcam)}{தஞ்சம் (tañcam)}]{\ta{தஞ்சம்} (tañcam) — refuge; shelter}

\label{lesson:TA-C303-tancam}



\begin{quote}
Before the new one: say the Tamil for a commotion, then the Tamil for a stampede.
\end{quote}

\subsection*{You'll want to know: \ta{தஞ்சம்}}
\textbf{\ta{தஞ்சம்}} — \emph{tañcam} — \textquotedblleft{}refuge; shelter\textquotedblright{}.

\begin{quote}
\ta{மக்கள்} \ta{பள்ளியில்} \ta{தஞ்சம்} \ta{புகுந்தார்கள்}.
\end{quote}

\begin{quote}
\emph{makkaḷ paḷḷiyil tañcam pukuntārkaḷ.}
\end{quote}

\begin{quote}
People took refuge in the school.
\end{quote}

\ta{தஞ்சம்} \ta{புகு}, to take refuge; \ta{தஞ்சம்} \ta{கொடு}, to give shelter.

\begin{cousinweb}
\textbf{\ta{தஞ்சம்}} (\emph{tañcam}), \textquotedblleft{}refuge, shelter\textquotedblright{}, is a Tamil word whose origin is not settled in the sources at hand.
\end{cousinweb}

\begin{grammarlens}[title={What you've built}]
One more word for telling what happened.
\end{grammarlens}

\subsection*{Your turn}
\begin{itemize}
\raggedright
  \item \emph{Say it:} \emph{tañcam}
  \item \emph{Say it:} \emph{tañcam}, once more
  \item \emph{Say it:} \emph{makkaḷ paḷḷiyil tañcam pukuntārkaḷ.} — People took refuge in the school.
  \item \emph{Recall:} say the Tamil for a commotion, then the Tamil for a stampede, then say \emph{tañcam} again
\end{itemize}

\begin{tcolorbox}[breakable,colback=teal!4,colframe=teal!35!black,title={Before you move on}]
What does \ta{தஞ்சம்} mean? (\textquotedblleft{}Refuge; shelter\textquotedblright{}.) Say it once more.
\end{tcolorbox}
\section[\texorpdfstring{\ta{விரட்டு} (viraṭṭu)}{விரட்டு (viraṭṭu)}]{\ta{விரட்டு} (viraṭṭu) — to drive away; to chase off}

\label{lesson:TA-C303-virattu}



\begin{quote}
Before the new one: say the Tamil for a stampede, then the Tamil for refuge.
\end{quote}

\subsection*{You'll want to know: \ta{விரட்டு}}
\textbf{\ta{விரட்டு}} — \emph{viraṭṭu} — \textquotedblleft{}to drive away; to chase off\textquotedblright{}.

\begin{quote}
\ta{காவலர்கள்} \ta{கூட்டத்தை} \ta{விரட்டினார்கள்}.
\end{quote}

\begin{quote}
\emph{kāvalarkaḷ kūṭṭattai viraṭṭiṉārkaḷ.}
\end{quote}

\begin{quote}
The police drove the crowd away.
\end{quote}

Past (I …): \ta{விரட்டினேன்}, \emph{viraṭṭiṉēṉ}.

What is chased off takes -\ta{ஐ}.

\begin{cousinweb}
\textbf{\ta{விரட்டு}} (\emph{viraṭṭu}) is native Dravidian (DEDR 5489), in a family of words for fear: to frighten off.
\end{cousinweb}

\begin{grammarlens}[title={What you've built}]
One more word for telling what happened.
\end{grammarlens}

\subsection*{Your turn}
\begin{itemize}
\raggedright
  \item \emph{Say it:} \emph{viraṭṭu}
  \item \emph{Say it:} \emph{viraṭṭu}, once more
  \item \emph{Say it:} \emph{kāvalarkaḷ kūṭṭattai viraṭṭiṉārkaḷ.} — The police drove the crowd away.
  \item \emph{Recall:} say the Tamil for a stampede, then the Tamil for refuge, then say \emph{viraṭṭu} again
\end{itemize}

\begin{tcolorbox}[breakable,colback=teal!4,colframe=teal!35!black,title={Before you move on}]
What does \ta{விரட்டு} mean? (\textquotedblleft{}To drive away; to chase off\textquotedblright{}.) Say it once more.
\end{tcolorbox}
